\begin{abstract}
Enterprise zkEVM L2 architectures in which each enterprise operates an
independent chain require separate zero-knowledge validity proofs to be
verified on the settlement layer. As the number of enterprises~$N$ grows,
aggregate L1 verification gas scales linearly, threatening economic
viability even on zero-fee networks. This paper evaluates three recursive
proof aggregation strategies---binary tree accumulation (halo2
snark-verifier), folding schemes (Nova/ProtoGalaxy with Groth16 decider),
and batch verification (SnarkPack)---for compressing~$N$ independent
halo2-KZG enterprise proofs into a single L1-verifiable proof. We
benchmark each strategy across $N \in \{1, 2, 4, 8, 16\}$ enterprises
with 30 iterations per configuration, measuring proof size, aggregation
time, and L1 verification gas. Folding with a Groth16 decider achieves
the strongest results: 15.3$\times$ gas reduction at $N{=}8$ (from
3.36M to 220K gas), constant 128-byte final proof size independent of~$N$,
and sub-linear aggregation time (11.75\,s at $N{=}8$, growing to only
13.75\,s at $N{=}16$). Binary tree accumulation delivers exactly
$N$-fold savings (8.0$\times$ at $N{=}8$) with production-proven tooling
but 180\,s aggregation time. SnarkPack provides 7.5$\times$ savings with
the fastest aggregation (200\,ms at $N{=}8$) but higher gas overhead.
All measurements reconcile within expected ranges against published
benchmarks from Polygon zkEVM, zkSync Era, Scroll, Nebra, and Axiom.
We identify four aggregation invariants---soundness, independence
preservation, order independence, and gas monotonicity---that must hold
in production deployments and recommend a phased adoption strategy:
binary tree accumulation for near-term deployment, transitioning to
ProtoGalaxy folding with Groth16 decider once the tooling matures.
\end{abstract}
